
\section[Unified model
]{Our unified model in a nutshell\label{sec:Single-Nutshell}}


In our non-disciplinary science-based model (the 'abstract physical neuron', as we call it), a neuron (highly simplified) is a sphere with two lipid layers (membrane), a similarly insulating output tube (axon), and an incompressible, partially ionized electrolyte fluid with different compositions and concentrations inside and outside.
They contain ions, simple chemical molecules, and complex biological formations.
These components  may have gradients and move at continuously changing speeds that can vary by several orders of magnitude during operation. 
The lipid layer is partially permeable (controlled or uncontrolled) for the components above (channels, mainly ion channels).

Between the membrane and the axon, there is a large amount of uncontrolled ion channels (\gls{AIS}). 
The elastic wall of the membrane is covered with ion layers, and therefore, a potential difference is created there (the resting potential).
The membrane represents \textit{a confined space in which ions behave differently from in an infinite space}; furthermore, it significantly alters the properties of the few-nanometer-thick layer near the membrane. In this layer, the concentration and potential change significantly with the distance from the membrane due to interactions between the membrane's ion layers and the electrolyte.
Ions sent by other neurons can enter the internal volume through controlled inputs (synapse).

%The membrane functions as a \hyperlink{ControlTheory}{\gls{PID} control circuit}.
Under a small external influence, a neuron remains in its resting state and regulates itself through non-controlled ion channels in its membrane.
If the external influence exceeds a threshold, the controlled channels suddenly release a large amount of ions into the internal space, putting the neuron into a transient state.
The membrane potential changes suddenly (within a few dozen nanoseconds), and, due to the mutual repulsion of ions, the pressure acting on the membrane wall and the ion concentration in the layer near the membrane also increase significantly and suddenly. At this point, the disciplines of electricity and thermodynamics can be combined: the pressure and the electric field are proportional to the number of ions in the volume.
The forces exerted on ions, calculated in a disciplinary way, are falsified 
in the sense that the other discipline also contributes inseparably; in a strict sense, neither of the two disciplines alone can describe ions' motion.
The detailed analysis shows that the elastic energy during the action potential
is about two orders of magnitude higher than the electrical energy, but the 
amount of ions is about three oders of magnitude lower than that of the non-dissociated molecules.
It that situation, the two effects are of the same order of magnitude,
so that they can modulate each other. The electric charge on the membrane's wall 
would produce a simple discharge current. However, its charge carriers are
delivered by the damped oscillation (a soliton) generated by the membrane.
The amplitude of the soliton changes the number of the charge carriers delivered,
so physiology measures an oscillating discharge current (that generates
a potential wave known as \gls{AP}). No significant mass transfer
takes place.

Since \textit{the same physical action generates the voltage gradient and the concentration gradient} (and, consequently, the pressure), they are inseparable and proportional to each other.
Generating a voltage change (clamping, direct current, magnetic pulse) or a pressure change (ultrasound, mechanical, shock waves) on the membrane mutually cause the other to change. Inputting ions (synaptic input) generates both changes.
The control circuit tries to restore the resting state, during which the mentioned changes decrease proportionally.
The transient state of the neuron can, in principle, be described by any of the mentioned quantities, but, as a practical matter, \textit{the effects of changes in the other entities must also be taken into account}: \textit{no single discipline, alone, can describe the changes}.
The elastic membrane implements an almost critically damped vibration, but the "vibrating mass" and the "spring force" are constantly changing.
In an electric resonant circuit, the capacitor's voltage and capacitance change due to the finite amount of charge.
\textit{The pressure wave caused by a force shock and the voltage wave caused by a voltage shock describe the same effect}.
In both cases, the other effect must be taken into account to some extent: an infinitesimal movement of the ion also changes both the pressure and the voltage.
Changes in pressure and charge alter the membrane's thickness (and therefore its capacity), and both voltage and pressure can cause a structural change (phase change, "melting") in the lipid chains that make up the membrane.
\index{melting!lipids}

Here, the difference comes to light: the carrier of the \gls{EM} interaction has no mass, while the thermodynamic one does.
In the case of the pressure wave, the speed is naturally finite (mass must be moved). For the electric wave, there are laws regarding only instantaneous interaction; so, they must be adapted to \textit{slow ion current}s based on physical assumptions.
\index{slow current}
It is simpler to describe the generation of the action potential as an electrical oscillation, and its propagation as a pressure wave (although this only represents the wave front well:
the \gls{AP} lasts as long as there is a non-equilibrium charge in the neuron: the repulsion between the excess ions pushes the ions out of the closed space, so the intensity of the wave continuously decreases).
Hyperpolarization is described in the electrical view as a capacitive current in the sequential oscillator circuit, in the thermodynamic view as a suction effect occurring in the opposite phase of the membrane. 
In both cases, however, it must be taken into account that charge and mass are inseparable in ions, and that their cross-disciplinary effects significantly reduce the applicability of the disciplinary laws.
Furthermore, as our force analysis suggests, the ions' motion results 
an interference between the longitudinal thermodynamic oscillation of elastic origin
and a discharge process of electrical origin. 
 
In simple words, a neuron combines different disciplines; it works like a two-tact biological internal combustion engine, and can be described theoretically as a special Carnot-type thermodynamic engine, see section~\ref{sec:Physics-Thermodynamics}. The synaptic input currents operate the ignition. The influx of $Na^+$ ions acts like an explosion, producing a large pressure impulse due to electrical repulsion between ions; see section~\ref{sec:Physics-Speculations}. The neuron's volume remains practically unchanged (representing an iso-volume process), and the pressure wave vibrates the particles (including dissociated ions) as a longitudinal wave, and the resulting offset potential moves the ions out of the neuron through the only available output channel: the axon. The pressure wave is similar to a sound wave, with negligible material transport. 
The force analysis shows that the pressure wave
causes a longitudinal vibration, while the changed potential adds
a potential-dependent longitudinal force component for moving the ions.

In the first stroke, the pressure impulse invests energy into compressing the elastic membrane and starting a compression wave (a soliton). In the second stroke, the elastic membrane pumps out some electrolyte (the driving force is a combined thermoelectrical/mechanical force rather than some magic battery) through the ion channels in the \gls{AIS}, which "resists" it, thereby generating a potential wave known as \gls{AP}; well described by electricity. The force analysis shows that the pressure wave
causes a longitudinal vibration, while the changed potential adds
a potential-dependent longitudinal force component for moving the ions. The compression
produces heat, as evidenced by a rise in temperature. The expansion consumes heat, as evidenced by a decrease in temperature. These are well-understood processes in thermodynamics.
However, they are still unexplained in biophysics, although they were observed
seven decades ago~\cite{HeatProductionNeuron:1958}.
The elastic membrane functions as a damped oscillator, which, in its negative amplitude phase, reverses the direction of electrolyte flow (and, consequently, the direction of the current carried by the ions).
In the discipline of electricity (when using the serial $RC$ model), it can be interpreted as a capacitive current, which in physiology is observed as hyperpolarization that inverts the polarity of the output voltage.
The ions represent both mass and charge simultaneously, and, correspondingly, one can observe, in addition to the electrical potential wave, mechanical, optical, and other features changing.
In different disciplines, pressure and voltage represent the same phenomenon, using concepts from various disciplines, although the 
speeds of those waves differ significantly.

 

Today, two major branches of disciplinary theories compete
for being applied to describe the neuronal operation \cite{CriticalElectricity:2018,ThermodynamicAPDrukarch:2022,PerspectivesNerveSignalPropagation:2024,ComparisonHHandSoliton:2010,HH_Potential_Controversies_2017,PiezoelectricNeuron:2025}. Those cited references also compare those theories.
The 2-decade-old Heimburg-Jackson thermodynamic theory is trying to break in alongside or displace the 7-decade-old all-electric Hodgkin-Huxley theory. The two theories apply one of the scientific disciplines developed for inanimate nature, unchanged, to living nature. There are unexplained observations for each of the two theories, and neither can describe the biological neuron's functioning without contradictions. 
In a disciplinary view, there is no chance of the two theories fusing. Although they can describe a wider range of phenomena together, the conditions under which they are applicable exclude their joint use. Moreover, to conceal contradictions and fill gaps, biophysics assumes (typically unspecified protein) mechanisms that contradict the fundamental principles of physics, and some of them even the disciplinary laws.




The electrical view essentially stems from the (Nobel Prize-winning) work of Hodgkin and Huxley \cite{HodgkinHuxley:1952}. They could not make a perfect job, mainly due to the lack of
discoveries made several decades later, including ours, about handling
the finite speed of the biological ionic currents (and, due to that, the dynamic features), which drastically change
their conclusions. In contrast with their \textit{empirical description}, which delivered mathematically formulated measured observations, our discussion, although it follows essentially
the same principles, reinterprets and pinpoints the used fundamental terms, sets up a physical model, and \textit{explains} the physically underpinned processes. ``However, the theory could not explain the physical
phenomena such as reversible heat changes, density changes,
and geometrical changes observed in the experiments'' \cite{ComparisonHHandSoliton:2010}.
The thermodynamic view roots in the (possibly Nobel-prize-winning)
idea of Heimburg and Jackson \cite{SolitonPropagation:2005}. They proposed that the action potential is essentially a sound wave (a soliton). "However, there are several other questions that this has to answer like ion flow involvement in nerve signal propagation as stated by the \gls{HH} model and also the faster propagation in myelinated nerves than in unmyelinated" \cite{ComparisonHHandSoliton:2010}.
If we consider that the ion flow means charge propagation in a membrane tube where the specific capacity depends on the thickness (of the myelin sheath) of the wall \cite{VeghUnifiedCrossdisciplinaryModel:2026}, the faster propagation is not mysterious anymore.
Those apparent contradictions are simply consequences of ions' non-disciplinary features.
By using the correct (non-disciplinary) discussion of the phenomena,
we can explain all those observations in more detail in section~\ref{sec:Physics-Unified}.
%Given that
%our model natively connects charge and mass of ions~\cite{VeghNon-ordinaryLaws:2025}, it has a good chance of explaining the missing phenomena.

The electrical view essentially stems from the (Nobel Prize-winning) work of Hodgkin and Huxley \cite{HodgkinHuxley:1952}. They could not make a perfect job, mainly due to the lack of
discoveries made several decades later, including ours, about handling
the finite speed of the biological ionic currents (and, due to that, the dynamic features), which drastically change
their conclusions. In contrast with their \textit{empirical description}, which delivered \textit{mathematically formulated measured observations}, our discussion, although it follows essentially
the same principles, reinterprets and pinpoints the used fundamental terms, sets up a physical model, and \textit{explains} the physically underpinned processes. "However, the theory could not explain the physical
phenomena such as reversible heat changes, density changes,
and geometrical changes observed in the experiments" \cite{ComparisonHHandSoliton:2010}. Given that
our (non-disciplinary) model natively connects charge and mass of ions~\cite{VeghNon-ordinaryLaws:2025}, it can explain the missing phenomena.
The thermodynamic view roots in the (possibly Nobel-prize-winning)
idea of Heimburg and Jackson \cite{SolitonPropagation:2005}. They proposed that the action potential is essentially a pressure wave (a soliton). ``However, there are several other questions that this has to answer like ion flow involvement in nerve signal propagation as stated by the \gls{HH} model and also the faster propagation in myelinated nerves than in unmyelinated'' \cite{ComparisonHHandSoliton:2010}.
If we consider that the presence of ions in the neuron's closed volume means an unbalanced force component exerting only on ions (but not on the vibrating neutral molecules)
and the ion flow means charge propagation in a membrane tube where the tube's specific capacity depends on the thickness of the axonal membrane's myelin sheath (aka the distance across the membrane surfaces), the faster propagation is not mysterious anymore.


The excellent textbook's statement that "the resting membrane potential results
from the separation of charge across the
cell membrane"~\cite{PrinciplesNeuralScience:2013} is only half the truth; we tell the second half in section~\ref{sec:Physics-MembraneElectricity}. 
We refute their statement that "%by discussing how 
resting ion channels
establish and maintain the resting potential"~\cite{PrinciplesNeuralScience:2013}, page 126.
The ion channels are passive players. The interplay of 
the lipid condenser and the electrolytes on both sides
of the permeable membrane establishes the resting potential, and ionic gradients maintain it.
The textbook separates the neuron membrane's states into resting and transient states. 
It introduces resting and transient ion channels,
but does not introduce that the different physical processes in those states need different models. 
As we show in section~\ref{sec:Physics-Unified}, \textit{two different physical mechanisms
operate in the resting and the transient states}, so two different models are needed for describing the same components.

A primary reason neurophysiology has been misguided is that it has in mind a static picture of the dynamic life of neurons.
It started from the static
inanimate anatomic microscopic pictures, continued
with the static clamping investigations (where
the negative feedback of an external electric circuit eliminates the natural gradients) and underpinning the
theory by analyzing frozen (inanimate) samples
by electron microscopes.
Those investigations revealed a tremendous amount of crucial details,
but they obscured the idea
that those samples are just snapshots of a permanently
changing system.
Neuroscience projected the mechanisms of the static resting state to the dynamic transient state and attempted
to understand the transient state 
as perturbations~\cite{PerturbationNeuralComputation:2002} to the  (entirely different) resting state.
Of course, that attempt required more and more ad hoc non-scientific assumptions.
We agree that "Understanding and predicting molecular responses
in single cells upon chemical, genetic or mechanical perturbations is a core question in biology."~\cite{PerturbationSingleCell:2023}
However, it is not possible without understanding that the transient state is a state on its own right,
with different physical mechanisms, rather than a perturbation of the entirely different resting state.

Due to the lack of the needed expertize in physics,
biophysics introduced the idea of the "whatnot"~\cite{Schrodinger:1992} protein mechanisms.
Claims such as "pumps [by using protein mechanisms] \dots transport ions \textit{against their
electrical and chemical gradients}"~\cite{PrinciplesNeuralScience:2013}, page 101, 
or "the flux of ions through ion channels is passive,
requiring \textit{no expenditure of metabolic energy} by the
channels."~\cite{PrinciplesNeuralScience:2013}, page 107, are as scientific as in the pre-Newtonian age,
there was a notion related to Aristotelian philosophy. This concept gave credence to the “theory” that the celestial spheres and, later on, the planets were pushed or moved by celestial movers, celestial intelligences, or angels against physical forces. Life science does not have its laws of motion, in the sense as Newton's Laws in physics, at least at cellular level. 


The book~\cite{PrinciplesNeuralScience:2013} asks the central questions, "How do ionic [i.e., electrical and chemical] gradients contribute to the resting membrane potential? What prevents
the ionic gradients from dissipating by diffusion of
ions across the membrane through the resting channels?"
However, it leaves them essentially open by giving only a qualitative answer, discussing only membrane permeability, without explaining how the resting potential is created.
%
We demonstrate in section~\ref{sec:Physics-MembraneElectricity} that classical methods of electricity enable us to calculate the membrane's potential resulting from charge separation and polarization, why and how its resting value is created, furthermore, how it is regulated.
We make some physically plausible assumptions to provide numerical figures.

Our results underpin that "the crucial system in biology isn't a molecule or a molecular class whatsoever, but the interface created by biomolecules in water"~\cite{LivingSystemPhysics:2021}. We add that mainly
physical processes at various speeds (instead of molecular classes such as proteins) and ionic gradients (instead of protein mechanisms) control the processes.
More precisely, inseparable thermodynamic and electrical processes set up and maintain the potential established by the electrical and thermodynamic backbone defined by the lipid/protein structure of the cell and the ionic solutions it contains. It is one of the frequent cases when 
one effect implements a functionality (the backbone closely defines the frames of the operation), and the other (in this case, combined thermodynamics and electricity) corrects it when it implements a complex operational (dynamical) functionality:
"stimulated phase transitions enable the phase-dependent processes to replace each other ... one process to build and the other
to correct"~\cite{BiologicalConservationLaw:2017}.

